% Part: many-valued-logic
% Chapter: three-valued-logics
% Section: kleene

\documentclass[../../../include/open-logic-section]{subfiles}

\begin{document}

\olfileid{mvl}{thr}{skl}

\olsection{د کليني منطقونه}

سټيفن کليني دوه درې ارزښته منطقونه وړاندې کړل۔ د هغو انګېزه داسې منطق
و چې د رښتياوالي قيمتونه د محاسبوي کړنلارو پايلې ګڼي: يوه کړنلاره
کېدے شي~$\True$ يا~$\False$ ورکړي، خو کېدے شي هېڅکله و نۀ درېږي۔
په دې حالت کښې اړوند د رښتياوالي قيمت ناټاکلے وي او د~$\Undef$
په قيمت ښودل کېږي۔

د يوې قضيې~$!A$ د نفي د محاسبې دپاره به لومړے د~$!A$ قيمت
محاسبه کړئ او بيا به د پايلې مخالف قيمت ورکړئ۔ که د~$!A$ محاسبه
پای ته و نۀ رسېږي، نو ټوله کړنلاره هم پای ته نۀ رسېږي؛ نو د~$\Undef$
نفي~$\Undef$ ده۔

د عطف~$!A \land !B$ د محاسبې دپاره دوه لارې شته: په يوه کښې
لومړے~$!A$ او بيا~$!B$ محاسبه کوو۔ که دواړو~$\True$ ورکړ، د عطف
پايله~$\True$ وي، او کنه~$\False$ وي۔ که له دغو دوو محاسبو يوه
هم و نۀ درېږي، ټوله کړنلاره پای ته نۀ رسېږي۔ نو په دې پرله‌پسې
تفسير کښې که د عطف يو اړخ ناټاکلے وي، ټول عطف هم ناټاکلے وي۔
د فصل په باب هم همداسې ده۔

خو که~$!A$ او~$!B$ په هممهاله ډول وارزوو، تر دې ښه پايله
تر لاسه کولے شو۔ که له دوو کړنلارو يوه ودرېږي او~$\False$
ورکړي، درېدلے شو، ځکه چې د عطف ځواب به خامخا دروغ وي۔ نو په
دې حالت کښې د عطف يو اړخ چې دروغ وي، ټول عطف دروغ وي، که څه هم
بل اړخ ناټاکلے وي۔ په همدې ډول، د فصل د هممهالې محاسبې پر مهال
چې د يوه اړخ کړنلاره ودرېږي او رښتيا ورکړي، درېدلے شو: فصل به
خامخا رښتيا وي۔ نو دلته د مرکبې دعوې پايله معلومولے شو، که څه
هم يو اړخ ئې ناټاکلے وي۔ په دې تفسير کښې~$\Undef$ د «ناټاکلي»
پر ځای د «نامعلوم» په معنا لوستل کېدے شي۔

له دغو دوو تفسيرونو څخه د کليني قوي او کمزوري منطقونه جوړېږي۔
شرطي د~$\lnot !A \lor !B$ په معادل ډول تعريفېږي۔

\begin{defn}
\emph{د کليني قوي منطق}~$\LogKs$ د لاندې ماتريس په وسيله تعريفېږي:
\begin{enumerate}
  \item معياري بياني ژبه~$\Lang L_0$ چې
  $\lnot$, $\land$, $\lor$, $\lif$ نښلوونکي لري۔
  \item د رښتياوالي د قيمتونو سټ~$V = \{\True, \Undef, \False\}$۔
  \item $\True$ يوازينے ټاکل شوے قيمت دے؛ يعنې~$V^+ = \{\True\}$۔
  \item د صدق تابعې په لاندې جدولونو کښې ورکړل شوې دي:
  \begin{center}
    \begin{tabular}{c|c} 
      $\tf{\lnot}$ & \\ 
      \hline  
      $\True$ & $\False$ \\ 
      $\Undef$ & $\Undef$ \\
      $\False$ & $\True$ 
    \end{tabular}
    \quad
    \begin{tabular}{c|ccc} 
      $\tf{\land}[\LogKs]$ & $\True$ & $\Undef$ & $\False$ \\ 
      \hline 
      $\True$ & $\True$ & $\Undef$ & $\False$ \\ 
      $\Undef$ & $\Undef$ & $\Undef$ & $\False$\\ 
      $\False$ & $\False$ & $\False$ & $\False$ 
    \end{tabular}
    \\[2ex]
    \begin{tabular}{c|ccc} 
      $\tf{\lor}[\LogKs]$ & $\True$ & $\Undef$ & $\False$ \\ 
      \hline 
      $\True$ & $\True$ & $\True$ & $\True$ \\ 
      $\Undef$ & $\True$ & $\Undef$ & $\Undef$ \\
      $\False$ & $\True$ & $\Undef$ & $\False$ 
    \end{tabular}
    \quad
    \begin{tabular}{c|ccc} 
      $\tf{\lif}[\LogKs]$ & $\True$ & $\Undef$ & $\False$ \\ 
      \hline 
      $\True$ & $\True$ & $\Undef$ & $\False$ \\ 
      $\Undef$ & $\True$ & $\Undef$ & $\Undef$  \\ 
      $\False$ & $\True$ & $\True$ & $\True$ 
    \end{tabular}
  \end{center} 
\end{enumerate}
\end{defn}

\begin{defn}
\emph{د کليني کمزورے منطق}~$\LogKw$ د لاندې ماتريس په وسيله تعريفېږي:
\begin{enumerate}
  \item معياري بياني ژبه~$\Lang L_0$ چې
  $\lnot$, $\land$, $\lor$, $\lif$ نښلوونکي لري۔
  \item د رښتياوالي د قيمتونو سټ~$V = \{\True, \Undef, \False\}$۔
  \item $\True$ يوازينے ټاکل شوے قيمت دے؛ يعنې~$V^+ = \{\True\}$۔
  \item د صدق تابعې په لاندې جدولونو کښې ورکړل شوې دي:
  \begin{center}
    \begin{tabular}{c|c} 
      $\tf{\lnot}$ & \\ 
      \hline  
      $\True$ & $\False$ \\ 
      $\Undef$ & $\Undef$ \\
      $\False$ & $\True$ 
    \end{tabular}
    \quad
    \begin{tabular}{c|ccc} 
      $\tf{\land}[\LogKw]$ & $\True$ & $\Undef$ & $\False$ \\ 
      \hline 
      $\True$ & $\True$ & $\Undef$ & $\False$ \\ 
      $\Undef$ & $\Undef$ & $\Undef$ & $\Undef$\\ 
      $\False$ & $\False$ & $\Undef$ & $\False$ 
    \end{tabular}
    \\[2ex]
    \begin{tabular}{c|ccc} 
      $\tf{\lor}[\LogKw]$ & $\True$ & $\Undef$ & $\False$ \\ 
      \hline 
      $\True$ & $\True$ & $\Undef$ & $\True$ \\ 
      $\Undef$ & $\Undef$ & $\Undef$ & $\Undef$ \\
      $\False$ & $\True$ & $\Undef$ & $\False$ 
    \end{tabular}
    \quad
    \begin{tabular}{c|ccc} 
      $\tf{\lif}[\LogKw]$ & $\True$ & $\Undef$ & $\False$ \\ 
      \hline 
      $\True$ & $\True$ & $\Undef$ & $\False$ \\ 
      $\Undef$ & $\Undef$ & $\Undef$ & $\Undef$  \\ 
      $\False$ & $\True$ & $\Undef$ & $\True$ 
    \end{tabular}
  \end{center} 
\end{enumerate}
\end{defn}

\begin{prop}
  په~$\LogKs$ او~$\LogKw$ کښې هېڅ تاوتولوژي نشته۔
\end{prop}

\begin{proof}
  که ټولو !!{propositional variable}s~$p$ ته
  $\pAssign v(p) = \Undef$ ورکړل شي، نو د هر فارمول~$!A$
  د رښتياوالي قيمت~$\pValue v(!A) = \Undef$ وي؛ ځکه چې
  \[
    \tf{\lnot}(\Undef) = \tf{\lor}(\Undef, \Undef) = \tf{\land}(\Undef,
  \Undef) = \tf{\lif}(\Undef, \Undef) = \Undef
  \]
  په دواړو منطقونو کښې همداسې ده۔ ځکه چې په~$\LogKs$ او
  $\LogKw$ دواړو کښې~$\Undef \notin V^+$ دے، نو په دې
  !!{valuation} کښې~$!A$ ټاکل شوے قيمت نۀ اخلي۔
\end{proof}

که څه هم د کليني په قوي او کمزوري منطقونو کښې تاوتولوژۍ نشته،
خو د هغو د استلزام اړيکې بې‌محتوا نۀ دي۔

\begin{prob}
  لاندې کومې اړيکې (الف) د کليني په قوي او (ب) د کليني په
  کمزوري منطق کښې صدق کوي؟ د هرې يوې دپاره د رښتياوالي جدول ورکړئ۔
  \begin{enumerate}
    \item $p, p \lif q \Entails q$
    \item $p \lor q, \lnot p \Entails q$
    \item $p \land q \Entails p$
    \item $p \Entails p \land p$
    \item $p \Entails p \lor q$
  \end{enumerate}
\end{prob}

دميتري بوچوار~$\Undef$ د «بې‌معنا» په توګه تعبير کړ او هڅه ئې وکړه
چې د دروغجن د پارادوکس په څېر پارادوکسونه پرې حل کړي: هغه وټاکله
چې پارادوکس لرونکې جملې به~$\Undef$ قيمت اخلي۔ ده داسې منطق
وړاندې کړ چې په اصل کښې د کليني کمزورے منطق دے، خو اضافي
نښلوونکي لري۔ له هغو څخه دوه «بهرنۍ نفي» او د «ناټاکلتيا»
عملګر دي:
\begin{center}
  \begin{tabular}{c|c} 
    $\tf{\sim}$ & \\ 
    \hline  
    $\True$ & $\False$ \\ 
    $\Undef$ & $\True$ \\
    $\False$ & $\True$ 
  \end{tabular}
\quad
  \begin{tabular}{c|c} 
    $\tf{+}$ & \\ 
    \hline  
    $\True$ & $\False$ \\ 
    $\Undef$ & $\True$ \\
    $\False$ & $\False$ 
  \end{tabular}
\end{center}

\begin{prob}
  آيا د بوچوار په منطق کښې~$\sim$ د~$\lnot$ او~$+$ په وسيله
  تعريفولے شئ؟ يعنې داسې !!a{formula} ومومئ چې يوازې
  !!{propositional
  variable}~$p$ پکښې وي، $\sim$ پکښې نۀ وي،
  او تل د~$\mathord{\sim}p$ په شان د رښتياوالي قيمت واخلي۔
  د خپل ځواب د تصديق دپاره د رښتياوالي جدول ورکړئ۔
\end{prob}


\end{document}
